\documentclass[11pt]{article}
\usepackage[utf8]{inputenc}
\usepackage{amsmath,amssymb}
\usepackage[margin=1in]{geometry}
\usepackage{hyperref}
\emergencystretch=3em

\title{Shifted monomial moments}
\author{Claude, with Daniel Murfet}
\date{2026-09-07}

\begin{document}
\maketitle
\sloppy

\begin{abstract}
First step of the continuous-amplitude milestone (Astra~\#16, route~C: monomial moments, polynomial
density, bounded mass). Dressing the bare normal moment integral with a monomial $u^\gamma$ shifts
the exponents $h\mapsto h+\gamma$. Normalising by the scale $N^{-\lambda}(\log N)^{|J|-1}$ of the
\emph{undressed} integral: if $\gamma$ vanishes on the minimal set $J$, the normalised moment converges
to $\Gamma(\lambda)\beta^{-\lambda}/(|J|-1)!\cdot\prod_{i\in J}\frac1{2k_i}\prod_{i\notin J}\frac1{h_i+\gamma_i+1-2k_i\lambda}$
(\texttt{shiftedMoment\_tendsto\_of\_face}); otherwise it tends to $0$
(\texttt{shiftedMoment\_tendsto\_zero}). These are the monomial moments of the face-supported
limiting functional. Zero \texttt{sorry}/\texttt{axiom}.
\end{abstract}

\section{The things formalised}
{\small
\begin{verbatim}
theorem tendsto_log_pow_div_pow (a b : ℕ) (hab : a < b) :
    Tendsto (fun N => log N ^ a / log N ^ b) atTop (𝓝 0)
theorem ratioExp_le_shift / ratioExp_shift_eq_of_zero / ratioExp_lt_shift / shift_min
theorem shiftedMoment_tendsto_of_face (d) (h k : Fin (d+1) → ℕ) (hk) (l β) (hl) (hβ)
    (hmin) (hatt) (γ : Fin (d+1) → ℕ) (hγ : ∀ i, ratioExp h k i = l → γ i = 0) :
    Tendsto (fun N => monomialBoxReal (d + 1) (fun i => h i + γ i) k β N
        / (N ^ (-l) * log N ^ (multCount (ratioExp h k) l - 1))) atTop
      (𝓝 (monomialMixedConst (fun i => h i + γ i) k l β))
theorem shiftedMoment_tendsto_zero ... (hγ : ∃ j, ratioExp h k j = l ∧ γ j ≠ 0) :
    Tendsto (fun N => monomialBoxReal (d + 1) (fun i => h i + γ i) k β N
        / (N ^ (-l) * log N ^ (multCount (ratioExp h k) l - 1))) atTop (𝓝 0)
\end{verbatim}
}

\section{Mathematics}
Shifting never decreases a ratio, strictly increases it exactly where $\gamma_i>0$. For a face shift
the minimum and its multiplicity are unchanged and the mixed theorem (unit~181) applies to $h+\gamma$.
Otherwise the multiplicity of $\lambda$ strictly drops: if it stays positive, the same exponent
appears with a lower log power and $(\log N)^{m'-1}/(\log N)^{m-1}\to0$; if it drops to zero, every
shifted ratio exceeds $\lambda$, the new minimum $\lambda'>\lambda$ is attained, and the quotient of
power--log leading terms with $\lambda<\lambda'$ tends to zero (\texttt{tendsto\_quotient\_zero\_of\_lt}).

\section{Paper-facing meaning}
For the dressed normal moment $\int u^h\eta(u)e^{-Nu^{2k}}$ these limits are the moments of the
leading functional: Taylor monomials supported on the nonminimal coordinates $I\setminus J$
contribute at leading order with the residue factors $1/(h_i+\gamma_i+1-2k_i\lambda)$, while any
monomial involving a minimal coordinate is of lower order. This is the mechanism behind Astra~\#16's
qualification of eq.~\texttt{thm\_leading\_coeff} for mixed ratios.

\section{Lean notes}
\begin{itemize}
\item \texttt{set m' := multCount ... with hm'} does not fold occurrences inside hypotheses created
\emph{later}; \texttt{rw [← hm'] at hT} after obtaining them.
\item \texttt{unfold} fails on a goal whose \texttt{set} variables hide the constant; \texttt{rw [hm]}
first, then \texttt{unfold}.
\item \texttt{rw [if\_neg H]} takes the negation as an argument; the remaining goal is the numeric
inequality, not \texttt{H}.
\item After \texttt{unfold ratioExp} the beta-redex \texttt{(fun i => h i + γ i) i} blocks
\texttt{rw}; use \texttt{simp only}.
\end{itemize}

\end{document}
